\section{Component catalogue}
\label{sec:catalogue}

Table~\ref{tab:catalogue} is generated by \texttt{docs/whitepaper/catalogue.py}
directly from the component registry: names come from the \texttt{@register}
decorators, and the grouping from the module each class lives in. It is
regenerated as a prerequisite of building this PDF, so it cannot list a component
that does not exist or omit one that does --- which is the same guarantee, from
the same table, that the config schema and the editor's palette get
(§\ref{sec:architecture}).

Each component's parameters, defaults and one-line purpose are in the published
documentation at \url{https://goncamateus.github.io/segmentator/}, in English and
Portuguese, and in the repository under \texttt{docs/en/stages.md}. The same
generator emits that longer form; it is left out here because a reader who needs
a parameter's default needs it next to a config they are editing, not in a PDF.

\begin{table}[h]
  \small
  \caption{The \NumStages{} stages by module family, and the \NumSources{}
    sources and \NumSinks{} sinks.}
  \label{tab:catalogue}
  \begin{tabular}{@{}l p{0.78\textwidth}@{}}
    \toprule
    Family & Components \\
    \midrule
    \inputrows{generated/catalogue-compact.tex}
    \bottomrule
  \end{tabular}
\end{table}
